\documentclass[11pt]{article}
\usepackage[a4paper,margin=1.1in]{geometry}
\usepackage{amsmath,amssymb,amsthm}
\usepackage[T1]{fontenc}
\usepackage{lmodern}
\emergencystretch=3em

\newtheorem{definition}{Definition}
\newtheorem*{goal}{Goal}

\newcommand{\At}{\mathbb A_t}
\newcommand{\Atev}{\mathbb A_t^{\mathrm{ev}}}
\newcommand{\Ft}{\mathbb F_t}
\newcommand{\Ftev}{\mathbb F_t^{\mathrm{ev}}}
\newcommand{\Nat}{\mathcal N_{\mathbb A_t}}
\newcommand{\Natev}{\mathcal N_{\mathbb A_t^{\mathrm{ev}}}}
\newcommand{\HC}{\operatorname{HC}}
\newcommand{\ev}{\operatorname{ev}}

% ============================================================================
% STATUS: PROPOSED, NOT SIGNED.  Drafted 2026-08-30 by hypervisorA on the
% Overseer's in-session order; REBUILT the same evening after a six-lens
% adversarial verification panel (46 findings, 10 BLOCKING) so that every
% definition below is the adopted manuscripts' own construction, verbatim
% where possible.  On the Overseer's signature this text is installed as
% SandboxA/informal/goalA.tex; the 07-18 reductive-group goal.tex is left
% untouched pending the flagship decision at root.
%
% PROVENANCE: this states, and states only, what the adopted pair proves --
%   SP-2026-036 v3  (flagship, sha256 2c9b5079...3103)
%   SP-2026-037 v2  (companion, sha256 721e7423...1f4d)
% typed adoption 2026-08-24 (Established-informal), hash-closed pair,
% deposits immutable at OID f972b120.  Clause provenance: (i) and the forward
% half of (ii) are the companion's Theorem 1.2; the reverse half of (ii) is
% the flagship's Verma-lattice characterization; (iii) is the flagship's
% Theorem GoalA (whose statement includes the intrinsic-centre identity as a
% CLAIM, as here); (iv) is the flagship's coefficient-even corollary.  An
% earlier independent proof of (iii) is SP-2026-010 (unused by the pair).
% U_{A_t} is defined by the Newton--PBW span (the flagship's presentation);
% the companion defines it intrinsically and PROVES the two agree -- that
% agreement is part of the pair's content, not an assumption here.  The
% root vectors are the manuscripts' INTERVAL RECURSION (not braid vectors:
% the braid normalization differs by the non-unit (v-v^{-1}) and would
% change the integral form).  Conventions: q = v^2, base ring A_t (the
% Overseer's 2026-08-29 default ruling); no square roots of the family
% parameters (the cover A_z is proof-technique only in both manuscripts).
% Panel record: informal/ACTIVITY.md 2026-08-30, workflow wf_eac43831-081.
% ============================================================================

\title{\textbf{goal A} --- the centre of the even hybrid family quantum group of $GL_n$}
\date{}

\begin{document}
\maketitle

Fix $n\ge 1$ and put
\[
  q=v^{2},\qquad
  \At=\mathbb Z[v,v^{-1}][t_1^{\pm1},\ldots,t_n^{\pm1}],\qquad
  \Ft=\mathbb Q(v)[t_1^{\pm1},\ldots,t_n^{\pm1}],
\]
\[
  \Atev=\mathbb Z[q,q^{-1}][t_1^{\pm1},\ldots,t_n^{\pm1}],\qquad
  \Ftev=\mathbb Q(q)[t_1^{\pm1},\ldots,t_n^{\pm1}],
\]
with $v$ an indeterminate and $t_1,\ldots,t_n$ algebraically independent.
Let $\varepsilon_1,\ldots,\varepsilon_n$ be the standard basis of
$\mathbb Z^{n}$, equipped with $(\varepsilon_a,\varepsilon_b)=\delta_{ab}$;
put $\alpha_i=\varepsilon_i-\varepsilon_{i+1}$ ($1\le i<n$) and
$\beta_{ab}=\varepsilon_a-\varepsilon_b$ ($1\le a<b\le n$), and let
$\Phi^{+}=\{\beta_{ab}:1\le a<b\le n\}\subset\mathbb Z^{n}$ be the set of
positive roots, listed from left to right in the order
\[
  (1,2),\,(1,3),\ldots,(1,n),\,(2,3),\ldots,(n-1,n).
\]
Arrays are written $\mathbf a=(a_{ab})_{a<b}\in\mathbb Z_{\ge0}^{\Phi^{+}}$,
with $\mathbf 0$ the zero array.  For $Z$ the centre of an algebra we write
$Z(-)$.  Quantum integers and factorials are
\[
  [r]_v=\frac{v^{r}-v^{-r}}{v-v^{-1}},\qquad
  [r]_v!=\prod_{j=1}^{r}[j]_v,\qquad
  [0]_v!=1 .
\]

\begin{definition}[Reduced double of family quantum $GL_n$]\label{def:utilde}
$\widetilde U_{\Ft}$ is the $\Ft$-algebra with invertible pairwise commuting
generators $K_a^{+},K_a^{-}$ ($1\le a\le n$) and generators $e_i,F_i$
($1\le i<n$), subject to
\[
  K_a^{+}K_a^{-}=t_a,\qquad
  K_a^{\pm}\,e_i\,(K_a^{\pm})^{-1}=v^{\pm(\varepsilon_a,\alpha_i)}e_i,\qquad
  K_a^{\pm}\,F_i\,(K_a^{\pm})^{-1}=v^{\mp(\varepsilon_a,\alpha_i)}F_i,
\]
\[
  e_iF_j-F_je_i=\delta_{ij}\bigl(K_i^{+}K_{i+1}^{-}-K_i^{-}K_{i+1}^{+}\bigr),
\]
\[
  e_ie_j=e_je_i,\qquad F_iF_j=F_jF_i\qquad(|i-j|>1),
\]
\[
  e_i^{2}e_j-(v+v^{-1})e_ie_je_i+e_je_i^{2}=0,\qquad
  F_i^{2}F_j-(v+v^{-1})F_iF_jF_i+F_jF_i^{2}=0\qquad(|i-j|=1).
\]
Define interval root vectors by $e_{a,a+1}=e_a$, $F_{a,a+1}=F_a$ and, for
$b>a+1$,
\[
  e_{ab}=\frac{e_{a,b-1}\,e_{b-1,b}-v^{-1}e_{b-1,b}\,e_{a,b-1}}{v-v^{-1}},
  \qquad
  F_{ab}=F_{b-1,b}\,F_{a,b-1}-v\,F_{a,b-1}\,F_{b-1,b}.
\]
For $\mathbf a\in\mathbb Z_{\ge0}^{\Phi^{+}}$ the divided and ordinary PBW
words are
\[
  F^{(\mathbf a)}=\prod_{a<b}^{\longleftarrow}\frac{F_{ab}^{\,a_{ab}}}{[a_{ab}]_v!},
  \qquad
  e^{\mathbf a}=\prod_{a<b}^{\longrightarrow}e_{ab}^{\,a_{ab}},
\]
where the arrows read the displayed root list backwards and forwards.  Put
$\nu(\mathbf a)=\sum_{a<b}a_{ab}\,\beta_{ab}\in\mathbb Z^{n}$, and for
$\mu=\sum_a\mu_a\varepsilon_a\in\mathbb Z^{n}$ put
$K^{+}_{\mu}=\prod_a (K_a^{+})^{\mu_a}$.
\end{definition}

\begin{definition}[Torally even part]\label{def:even}
$\widetilde U_{\Ft}$ carries the $\mathbb Z^{n}/2\mathbb Z^{n}$-grading
determined on generators by
\[
  \deg K_a^{+}=\varepsilon_a,\qquad
  \deg K_a^{-}=-\varepsilon_a,\qquad
  \deg K^{+}_{\mu}=\mu,\qquad
  \deg e_i=0,\qquad
  \deg F_i=-\alpha_i,\qquad
  \deg t_a=0;
\]
every defining relation is homogeneous modulo $2\mathbb Z^{n}$ (in the mixed
relation both toral terms have degree $-\alpha_i\equiv\alpha_i$, the degree
of its left side), so the grading is well defined, and the degree-zero part
is a subalgebra.  The interval vectors are homogeneous with
$\deg e_{ab}=0$ and $\deg F_{ab}=-\beta_{ab}$, so
$\deg F^{(\mathbf a)}=-\nu(\mathbf a)$.  $U_{\Ft}$ is the degree-zero part
of $\widetilde U_{\Ft}$.
\end{definition}

\begin{definition}[Newton lattice; the hybrid form]\label{def:newton}
Put
\[
  Y_a=q^{\,n-a}\,t_a^{-1}\,(K_a^{+})^{2}\in\widetilde U_{\Ft},
  \qquad\text{equivalently}\qquad
  (K_a^{+})^{2}=q^{\,a-n}\,t_a\,Y_a .
\]
Granting the basis in clause (i) below, the $Y_a$ are algebraically
independent invertible elements of $U_{\Ft}$, so
$\Ft[Y_1^{\pm1},\ldots,Y_n^{\pm1}]$ is a Laurent polynomial ring; for
$\lambda\in\mathbb Z^{n}$ let
\[
  \ev_\lambda:\Ft[Y_1^{\pm1},\ldots,Y_n^{\pm1}]\longrightarrow\Ft,
  \qquad
  \ev_\lambda(Y_a)=q^{\lambda_a+n-a}\,t_a^{-1},
\]
be the $\Ft$-algebra map.  The \emph{Newton lattice} is
\[
  \Nat=\bigl\{\,f\in\Ft[Y_1^{\pm1},\ldots,Y_n^{\pm1}]:
  \ev_\lambda(f)\in\At\ \text{for every }\lambda\in\mathbb Z^{n}\,\bigr\},
\]
and the \emph{even hybrid form} is the $\At$-submodule
\[
  U_{\At}\;=\;\sum_{\mathbf a,\mathbf b\in\mathbb Z_{\ge0}^{\Phi^{+}}}
  F^{(\mathbf a)}\,K^{+}_{\nu(\mathbf a)}\,\Nat\,e^{\mathbf b}
  \;\subset\;U_{\Ft},
\]
each summand torally even since
$\deg\bigl(F^{(\mathbf a)}K^{+}_{\nu(\mathbf a)}\bigr)=0$; only finite sums
occur.  The symmetric group $S_n$ fixes $\Ft$ and permutes the variables
$Y_1,\ldots,Y_n$; the invariant Newton lattice is
$(\Nat)^{S_n}=\{f\in\Nat: wf=f\ \text{for every }w\in S_n\}$.
\end{definition}

\begin{definition}[Integral Verma lattices; Harish--Chandra
projection]\label{def:verma-hc}
Let $\widetilde U_{\Ft}^{\ge0}\subset\widetilde U_{\Ft}$ be the subalgebra
generated by all $K_a^{\pm}$ and all $e_i$.  For $\lambda\in\mathbb Z^{n}$
let $\chi_\lambda:\widetilde U_{\Ft}^{\ge0}\to\Ft$ be the $\Ft$-algebra map
with
\[
  \chi_\lambda(K_a^{+})=v^{\lambda_a},\qquad
  \chi_\lambda(K_a^{-})=t_a\,v^{-\lambda_a},\qquad
  \chi_\lambda(e_i)=0,
\]
writing $\Ft^{\chi_\lambda}$ for $\Ft$ viewed as a
$\widetilde U_{\Ft}^{\ge0}$-module through $\chi_\lambda$, let
$M_{\Ft}(\lambda)
=\widetilde U_{\Ft}\otimes_{\smash{\widetilde U_{\Ft}^{\ge0}}}\Ft^{\chi_\lambda}$
be the induced module, with highest vector $m_\lambda=1\otimes1$; its
integral lattice is
\[
  M_{\At}(\lambda)\;=\;\sum_{\mathbf a\in\mathbb Z_{\ge0}^{\Phi^{+}}}
  \At\,F^{(\mathbf a)}\,K^{+}_{\nu(\mathbf a)}\,m_\lambda .
\]
Granting the basis in clause (i) below, every element of $U_{\Ft}$ has a
unique PBW coefficient on the row $\mathbf a=\mathbf b=\mathbf 0$, and that
coefficient lies in $\Ft[Y_1^{\pm1},\ldots,Y_n^{\pm1}]$; the
\emph{Harish--Chandra projection} is the $\Ft$-linear map
\[
  \HC:U_{\Ft}\longrightarrow\Ft[Y_1^{\pm1},\ldots,Y_n^{\pm1}]
\]
retaining that coefficient.
\end{definition}

\begin{goal}\begin{enumerate}
  \item[(i)] The PBW sums are direct: the monomials
  $F^{(\mathbf a)}(K_1^{+})^{\eta_1}\cdots(K_n^{+})^{\eta_n}e^{\mathbf b}$,
  for $\mathbf a,\mathbf b\in\mathbb Z_{\ge0}^{\Phi^{+}}$ and
  $\eta\in\mathbb Z^{n}$, form an $\Ft$-basis of $\widetilde U_{\Ft}$; the
  degree-zero part decomposes as
  \[
    U_{\Ft}\;=\;\bigoplus_{\mathbf a,\mathbf b\in\mathbb Z_{\ge0}^{\Phi^{+}}}
    F^{(\mathbf a)}\,K^{+}_{\nu(\mathbf a)}\,
    \Ft[Y_1^{\pm1},\ldots,Y_n^{\pm1}]\,e^{\mathbf b};
  \]
  the sum defining $U_{\At}$ is direct, $U_{\At}$ is an $\At$-subalgebra of
  $U_{\Ft}$ with $\Ft\otimes_{\At}U_{\At}\simeq U_{\Ft}$, and each sum
  defining $M_{\At}(\lambda)$ is direct.
  \item[(ii)] The even hybrid form is the common preserver of the integral
  Verma lattices:
  \[
    U_{\At}\;=\;\bigl\{\,x\in U_{\Ft}:
    x\,M_{\At}(\lambda)\subseteq M_{\At}(\lambda)\ \text{for every }
    \lambda\in\mathbb Z^{n}\,\bigr\}.
  \]
  \item[(iii)] The intrinsic centre satisfies
  $Z(U_{\At})=Z(\widetilde U_{\Ft})\cap U_{\At}$, and Harish--Chandra
  projection restricts to an isomorphism of $\At$-algebras
  \[
    \HC:\;Z\bigl(U_{\At}\bigr)\;\xrightarrow{\ \sim\ }\;
    \bigl(\Nat\bigr)^{S_n}.
  \]
  \item[(iv)] With
  \[
    \Natev=\bigl\{\,f\in\Ftev[Y_1^{\pm1},\ldots,Y_n^{\pm1}]:
    \ev_\lambda(f)\in\Atev\ \text{for every }\lambda\in\mathbb Z^{n}\,\bigr\}
  \]
  and $\mathfrak Z_{\Atev}=\{C\in Z(U_{\At}):\HC(C)\in(\Natev)^{S_n}\}$ the
  coefficient-even central form, $\HC$ restricts to an isomorphism of
  $\Atev$-algebras
  $\mathfrak Z_{\Atev}\xrightarrow{\ \sim\ }(\Natev)^{S_n}$, and
  multiplication induces
  $\At\otimes_{\Atev}\mathfrak Z_{\Atev}\xrightarrow{\ \sim\ }Z(U_{\At})$.
\end{enumerate}
\end{goal}

\end{document}
